%% ============================================================
%% Statistical Reframing Notes
%% These are revision notes for rewriting Results sections, not a standalone section.
%% The key changes to make throughout the paper:
%% ============================================================

%% PROBLEM (from R3):
%% prerequisite_chain rho=0.686 [.576,.775] overlaps with teacher rho=0.555 [.422,.664]
%% We cannot claim individual prompt X beats individual prompt Y.
%%
%% HONEST FRAMING:
%% The finding is GROUP-level: item-analysis prompts as a class outperform
%% population-only prompts as a class. Within groups, individual prompts
%% are not reliably distinguishable.
%%
%% CHANGES NEEDED:
%%
%% 1. Abstract: Change "Prompts directing the model to enumerate structural
%%    features achieve rho=0.65--0.69" to frame as group-level finding.
%%    Keep the range but add "as a class" or similar qualifier.
%%
%% 2. Section 5.1 title: "Item Analysis Prompts Achieve rho=0.65--0.69"
%%    is fine as a group-level claim (it's about the class of prompts).
%%    But the text should note that individual prompts within this group
%%    are not reliably distinguishable from each other.
%%
%% 3. Table 5 (screening): Add a note that CIs overlap substantially
%%    between adjacent ranks. The ranking is a point estimate ordering,
%%    not a statistically significant ordering.
%%
%% 4. Contributions paragraph: Reframe (2) from individual prompt claims
%%    to the group-level distinction (item-centric vs population-centric).
%%
%% 5. Table 6 (generalization): R2's concern about "best temperature per
%%    dataset" is valid. Add a column showing fixed t=1.0 results alongside.
%%    The main claims should use fixed temperature; best-temperature results
%%    are supplementary showing the ceiling.
%%
%% 6. Add explicit paragraph in Results or Method:
%%    "With 140 items, the 95% CI for rho=0.69 spans [.58, .78]---a width
%%    of 0.20. Individual prompt rankings should be interpreted as point
%%    estimates, not as evidence of reliable differences between similarly-
%%    performing prompts. Our primary claims are at the group level:
%%    item-analysis prompts as a class vs. population-only prompts as a class."
%%
%% 7. Discussion "Methodological Recommendations": Add sample size guidance
%%    for distinguishing individual prompts. With 140 items, two prompts
%%    must differ by ~0.15 rho to be distinguishable at p<.05.

%% TEMPERATURE REFRAMING:
%%
%% Current Table 6 uses "best temperature per dataset" which R2 rightly
%% flags as overfitting. For a real application, you don't have ground truth
%% to select the best temperature.
%%
%% Fix: Primary results at fixed t=1.0 (the default). Show best-temperature
%% results in a supplementary column labeled "oracle temperature" to make
%% clear this is an upper bound, not a practical recommendation.
%%
%% For the BEA comparison specifically: our calibrated results used the
%% training set for calibration, which is legitimate. But temperature
%% selection on the test set is not. Report the t=1.0 BEA result as primary.

%% FEW-SHOT PROMPTING (R2):
%%
%% R2 asks why few-shot is omitted from the design space. Honest answer:
%% (1) Few-shot requires exemplar items with known difficulty, which
%%     contradicts the zero-shot use case we target (new items without
%%     calibration data).
%% (2) Exemplar selection introduces a confound: which items you show
%%     as examples affects the estimate, creating a new design dimension
%%     we could not systematically explore within budget.
%% (3) Practical note: with long item texts + rubrics, few-shot prompts
%%     approach context limits and significantly increase per-item cost.
%%
%% Add a paragraph in Method or Discussion acknowledging this gap and
%% the reasons for the design choice.
